\subsubsection{Recognition and synthesis of one-owner tables}
\label{app:owner-synthesis}

Guard comparison decides equality of given terms. We now characterize
which abstract tables are terms. Fix input positions $N$, $|N|=n$, and
a candidate root $r$. A canonical candidate $\mathcal T$ has
$\mathcal T_r=\{(\varnothing,\varnothing)\}$. For $j\ne r$ its
family $\mathcal T_j$ is an antichain of pairs $(A,B)$ under coverage,
with
\[
 r\in A,\qquad j\notin A,\qquad B\cap(A\cup\{j\})=\varnothing.
\]
Empty families are allowed. These conditions are necessary by root
status and the canonical-table argument above. Write
\[
 \mathsf{ok}(j;C,D)\quad\Longleftrightarrow\quad
 \exists(A,B)\in\mathcal T_j:\quad
 A\subseteq C,\quad B\subseteq C\cup D\cup\{j\}.
\]
Thus the contribution with source $j$ and tests $(C,D)$ is already
permitted by the candidate. This predicate is monotone as tests grow,
where an ambiguity test implies a nonowner test.

For each nonroot guard $(j,A,B)$ run the following finite saturation.
Set $U=A\cup B\cup\{j\}$ and initialize $C=\{r\}$, $D=\varnothing$.
Repeatedly choose $w\in A$ and $z\in U$ for which, putting
$E=D\cup\{w,z\}$,
\begin{equation}\label{owner:eq:saturation}
 \mathsf{ok}(w;C,E)\quad\text{and}\quad
 \mathsf{ok}(z;C,E),
\end{equation}
and replace
\[
 (C,D)\longmapsto
 \bigl(C\cup\{w\},\ E\setminus(C\cup\{w\})\bigr)
\]
whenever this changes the state. At a fixed point accept the guard if
\begin{equation}\label{owner:eq:finish}
 \mathsf{ok}(j;C,D),\quad\text{or}\quad
 \exists z\in U:\ 
 \mathsf{ok}(j;C,D\cup\{j,z\})\ \land\
 \mathsf{ok}(z;C,D\cup\{j,z\}).
\end{equation}
The root guard is accepted without a test. Reject a noncanonical input
or any table with a rejected guard.

\begin{theorem}[One-owner table synthesis]\label{owner:thm:synthesis}
A canonical table is realized by a term in $\{\del,H_a\}$ if and only
if every guard passes this saturation test. Recognition is polynomial
in the table description, and an accepted table with $M$ guards has
a realizing native tree of size $O(nM)$, constructible in polynomial
time. The outcome does not depend on the order of enabled updates.
\end{theorem}
\begin{proof}
Throughout saturation, $C\subseteq A$ and $C\cup D\subseteq U$.
The rank $2|C|+|D|$ strictly increases at a changing update and is at
most $2n$. An update never removes a test: replacing $N_w$ by $M_w$
strengthens it. Enabled updates remain enabled under strengthening.
Consequently exhaustive application terminates at the least fixed point
above the initial state, independently of order.

We first construct a term for an accepted guard. A one-hole path starts
with $[-]\del x_r$, so the tests above the hole are
$(\{r\},\varnothing)$ and its required root is $r$.
An update using $w,z$ replaces the hole by
\[
 H_a(x_r,\ [-]\del x_w,\ x_z).
\]
The new hole again has root $r$. The two new contributions with sources
$w,z$ have tests $(C,D\cup\{w,z\})$; they are permitted precisely
by \eqref{owner:eq:saturation}. The contribution of $x_r$ is always
permitted. The tests above the new hole are exactly the updated state:
the $H_a$ frame adds $M_r,N_w,N_z$, and the reset frame adds $M_w$.
Here $r\in C$ already.

For the first alternative of \eqref{owner:eq:finish}, fill the hole with
$x_j\del x_r$. For the second, fill it with $H_a(x_r,x_j,x_z)$.
Every contribution is permitted by $\mathcal T$, and the displayed
occurrence of $j$ has tests implied by $(A,B)$: only tests in $A$ and
$U$ were introduced, and $N_j$ follows from the source itself.
Thus this path is contained in the candidate and covers the requested
guard. It has root $r$ and $O(n)$ nodes. Join these paths by reset.
Same-root reset is their union, including when the common root is
atomic, so the resulting term has exactly the candidate table.
There are $M$ paths, giving size $O(nM)$.

For necessity, take a term with table $\mathcal T$ and an occurrence
supplying a minimal nonroot guard $(j,A,B)$. Follow its route from the
root. At a reset, descending into its right child adds no tests and
retains the root; discard the left child. Descending into the left child
retains a frame $[-]\del x_v$, where $v$ is the right child's root.
At $H_a$, descending into its first child adds no tests and retains the
root; discard the other children. Descending into either other child
retains a frame $H_a(x_v,[-],x_z)$, using symmetry of the last two
arguments. Replace each off-path child by its root variable.
This pruning preserves the chosen occurrence's tests and only removes
other contributions, since each child's root contribution was already
present. The resulting path is therefore contained in $\mathcal T$.

Each nonterminal root $v$ of this path belongs to $A$, since leaving
it adds $M_v$. Every companion $z$ belongs to $U$, since it adds $N_z$.
Compare the path from outside to inside with a saturated state $(C,D)$.
Inductively its accumulated tests are implied by $(C,D)$, and its
current nonterminal root lies in $C$. This holds initially for $r$.
If a reset frame leads to another nonterminal root $w$, the path's
contribution of $w$ is permitted. Thus
$\mathsf{ok}(w;C,D\cup\{w\})$ holds. The saturation rule with
$w=z$ forces $w\in C$, and the induction continues. If an $H_a$
frame leads to a nonterminal root $w$ with companion $z$, both of
their contributions are permitted under $D\cup\{w,z\}$.
At a fixed point the rule forces $w\in C$ and $z\in C\cup D$,
again giving the induction. The last frame leads to the source $j$.
A reset gives the first alternative of \eqref{owner:eq:finish}, and
an $H_a$ frame gives its second alternative. The guard is accepted.

Finally, each guard needs at most $2n$ changing updates. A full scan
tests at most $n^2$ pairs using coverage against the supplied table.
Both recognition and recording the constructing frames are polynomial.
\end{proof}

This also describes the one-owner term clone, including substitution.
If inner tables have roots $r_i$, substitute them into an outer table
by replacing each outer guard $(i,A,B)$ and inner guard $(j,C,D)$ of
the $i$th table by
\[
 \left(j,\ C\cup\{r_\ell:\ell\in A\},\
             D\cup\{r_\ell:\ell\in B\}\right),
\]
then normalize. The composite root is $r_{r}$, where $r$ is the outer
root. Root status and the membership recurrence prove this formula
directly. The synthesis theorem identifies exactly the tables on which
these operations act; it is stronger than enumerating closure under
the original operations.

Recognition makes the free algebras computable. Through arity three
they have $1,4,27$ elements, which agrees with complete four-state
term-function closure. At arity four, $201$ of the $8000$ canonical
candidates with a fixed root are realizable, so there are $804$
four-ary term functions. These counts are computed by
\path{checks/check_owner_synthesis.py}, which confirms each candidate
by independent exhaustive path search.
